\documentclass[../../main.tex]{subfiles}
\begin{document}

{\bf [解答]}

% TODO: 四面体と球
\begin{figure}[htb]
  \centering
  % 俯角 65°，方位角 115°
  \tdplotsetmaincoords{65}{60}
  \begin{tikzpicture}[tdplot_main_coords, scale=2.5, line join=round, >=stealth, every node/.style={font=\small}]

    % --- 1. 厳密な頂点座標 ---
    % B: 原点 (0, 0, 0)
    % C: (2, 0, 0)
    % D: (1, sqrt(3), 0) ≈ (1, 1.7321, 0)
    % E: CD の中点 (1.5, sqrt(3)/2, 0) ≈ (1.5, 0.8660, 0)
    % A: (0.75, sqrt(3)/4, 1.5) ≈ (0.75, 0.4330, 1.5)
    \coordinate (B) at (0, 0, 0);
    \coordinate (C) at (2.0, 0, 0);
    \coordinate (D) at (1.0, 1.732051, 0);
    \coordinate (E) at (1.5, 0.866025, 0);
    \coordinate (A) at (0.75, 0.433013, 1.5);

    % --- 3. 四面体の面（透過塗りつぶし）---
    % 底面 BCD (xy 平面上)
    \fill[gray!10, opacity=0.6] (B) -- (C) -- (D) -- cycle;
    % 側面 ACD (手前右)
    \fill[blue!10, opacity=0.5] (A) -- (C) -- (D) -- cycle;
    % 側面 ABC (手前左)
    \fill[teal!12, opacity=0.5] (A) -- (B) -- (C) -- cycle;

    % --- 4. 稜線の描画（奥側の隠れ線を破線に）---
    % 辺 BD（奥側）
    \draw[dashed, gray!60, thick] (B) -- (D);

    % 手前の外枠稜線（実線）
    \draw[very thick, black] (B) -- (C);
    \draw[very thick, black] (C) -- (D);
    \draw[very thick, black] (A) -- (B);
    \draw[very thick, black] (A) -- (C);
    \draw[very thick, black] (A) -- (D);

    % 対称面の補助線 AE, BE (破線)
    \draw[dashed, teal!80!black, thick] (A) -- (E);
    \draw[dashed, teal!80!black, thick] (B) -- (E);

    % --- 5. 頂点・中点プロットとラベル ---
    \node[circle, fill, inner sep=1.2pt] at (B) {}; \node[below left=2pt, font=\normalsize\bfseries] at (B) {$B(O)$};
    \node[circle, fill, inner sep=1.2pt] at (C) {}; \node[below right=2pt, font=\normalsize\bfseries] at (C) {$C$};
    \node[circle, fill, inner sep=1.2pt] at (D) {}; \node[above right=1pt, font=\normalsize\bfseries] at (D) {$D$};
    \node[circle, fill, inner sep=1.2pt] at (A) {}; \node[above=2pt, font=\normalsize\bfseries] at (A) {$A$};

    % 中点 E のプロット
    \node[circle, fill=teal!85!black, inner sep=1.0pt] at (E) {}; \node[below right=1pt, font=\scriptsize\bfseries, teal!90!black] at (E) {$E$};

    % --- 6. 各辺の長さラベル ---
    \node[above left=1pt, font=\footnotesize]  at ($ (A)!0.5!(B) $) {$\sqrt{3}$};
    \node[below=2pt, font=\footnotesize]       at ($ (B)!0.5!(C) $) {$2$};
    \node[below right=1pt, font=\footnotesize] at ($ (C)!0.5!(D) $) {$2$};
    \node[above=2pt, font=\footnotesize, gray!70!black] at ($ (B)!0.6!(D) $) {$2$};
    \node[right=2pt, font=\footnotesize]       at ($ (A)!0.5!(C) $) {$2$};
    \node[above right=1pt, font=\footnotesize] at ($ (A)!0.5!(D) $) {$2$};

  \end{tikzpicture}
  \caption{底面 $\triangle BCD$ を $xy$ 平面上に配置した四面体 $ABCD$ の 3D 見取り図}
  \label{fig:tetrahedron_3d_exact_coordinates}
\end{figure}

題意の球の中心を $O$ とする．$CD$ の中点を $E$ とし $\triangle ABE$ で切断すると下図のようになる．

% triangle ABE
\begin{figure}[htb]
  \centering
  \begin{tikzpicture}[scale=2.8, >=stealth, every node/.style={font=\small}]

    % --- 1. 座標定義 (底辺 BE, 頂点 A の二等辺/正三角形 ABE) ---
    % AB = sqrt(3), AE = sqrt(3), BE = sqrt(3)
    % 幅 W = sqrt(3) ≈ 1.732, 高さ H = 1.5
    % B(-W/2, 0), E(W/2, 0), A(0, 1.5)
    \def\wHalf{0.866025} % sqrt(3)/2
    \def\hVal{1.5}       % 3/2

    \coordinate (A) at (0, \hVal);
    \coordinate (B) at (-\wHalf, 0);
    \coordinate (E) at (\wHalf, 0);

    % 底辺 BE 上の点 Q: BQ = 2/3 * sqrt(3), QE = 1/3 * sqrt(3)
    % B から E への比率 2 : 1
    \coordinate (Q) at ($ (B)!{2.0/3.0}!(E) $);

    % 辺 AE 上の点 P: AP = 2/3 * sqrt(3), PE = 1/3 * sqrt(3)
    % A から E への比率 2 : 1
    \coordinate (P) at ($ (A)!{2.0/3.0}!(E) $);

    % 辺 AB の中点 R
    \coordinate (R) at ($ (A)!0.5!(B) $);

    % 球心 O: Q を通り BE ⊥ な直線 (x = x_Q) と P を通り AE ⊥ な直線の交点
    % Q = (wHalf/3, 0) = (sqrt(3)/6, 0)
    % AE の傾き = -1.5 / (sqrt(3)/2) = -sqrt(3)
    % OP ⊥ AE より OP の傾き = 1/sqrt(3)
    % OP の方程式: y - y_P = (1/sqrt(3)) * (x - x_P)
    % x = x_Q = sqrt(3)/6 を代入して y_O を厳密算出
    \path let \p1 = (P), \p2 = (Q) in
    coordinate (O) at (\x2, {\y1 + (1.0/sqrt(3.0))*(\x2 - \x1)});

    % --- 2. 三角形 ABE の枠線 ---
    \draw[very thick, black] (A) -- (B) -- (E) -- cycle;

    % --- 3. 垂線と内部補助線 ---
    % 垂線 OQ ⊥ BE
    \draw[very thick, black] (O) -- (Q);
    \node[below=2pt, font=\normalsize\bfseries] at (Q) {$Q$};

    % 垂線 OP ⊥ AE
    \draw[very thick, black] (O) -- (P);
    \node[above right=1pt, font=\normalsize\bfseries] at (P) {$P$};

    % 垂線 OR ⊥ AB (手書きスケッチの破線)
    \draw[dashed, black, thick] (O) -- (R);
    \node[above left=1pt, font=\normalsize\bfseries] at (R) {$R$};

    % 線分 OE
    \draw[very thick, black] (O) -- (E);

    % 線分 OA (長さ r)
    \draw[very thick, black] (O) -- (A);
    \node[above left=1pt, font=\normalsize\bfseries] at ($ (O)!0.5!(A) $) {$r$};

    % --- 4. 直角記号 ---
    % Q における直角 (OQ ⊥ BE)
    \draw[thin, black]
    ($ (Q) + (0.07, 0) $) --
    ($ (Q) + (0.07, 0.07) $) --
    ($ (Q) + (0, 0.07) $);

    % P における直角 (OP ⊥ AE)
    \pgfmathsetmacro{\angAE}{atan2(-\hVal, \wHalf)} % ≈ -60°
    \draw[thin, black]
    ($ (P) + ({\angAE}:-0.07) $) --
    ($ (P) + ({\angAE}:-0.07) + ({\angAE-90}:0.07) $) --
    ($ (P) + ({\angAE-90}:0.07) $);

    % R における直角 (OR ⊥ AB)
    \pgfmathsetmacro{\angAB}{atan2(-\hVal, -\wHalf)} % ≈ -120°
    \draw[thin, black]
    ($ (R) + ({\angAB}:-0.07) $) --
    ($ (R) + ({\angAB}:-0.07) + ({\angAB+90}:0.07) $) --
    ($ (R) + ({\angAB+90}:0.07) $);

    % --- 5. 頂点プロットとラベル ---
    \node[circle, fill, inner sep=1.0pt] at (A) {}; \node[above=2pt, font=\normalsize\bfseries] at (A) {$A$};
    \node[circle, fill, inner sep=1.0pt] at (B) {}; \node[below left=2pt, font=\normalsize\bfseries] at (B) {$B$};
    \node[circle, fill, inner sep=1.0pt] at (E) {}; \node[below right=2pt, font=\normalsize\bfseries] at (E) {$E$};
    \node[circle, fill, inner sep=1.0pt] at (O) {}; \node[left=2pt, font=\normalsize\bfseries] at (O) {$O$};
    \node[circle, fill, inner sep=0.8pt] at (P) {};
    \node[circle, fill, inner sep=0.8pt] at (Q) {};
    \node[circle, fill, inner sep=0.8pt] at (R) {};

    % --- 6. 寸法ラベル（手書きスケッチの引き出し線・配置を再現）---
    % 辺 AB の長さ sqrt(3)
    \node[left=6pt, font=\footnotesize] at (R) {$\sqrt{3}$};

    % 辺 AE の分割寸法 (2/3 * sqrt(3), 1/3 * sqrt(3))
    \draw[thick] ($ ($ (A)!0.5!(P) $) + (-0.05, -0.03) $) -- ($ ($ (A)!0.5!(P) $) + (0.05, 0.03) $);
    \draw[thick] ($ ($ (A)!0.5!(P) $) + (-0.05, -0.01) $) -- ($ ($ (A)!0.5!(P) $) + (0.05, 0.05) $);
    \node[above right=2pt, font=\footnotesize] at ($ (A)!0.5!(P) $) {$\frac{2}{3}\sqrt{3}$};
    \node[above right=2pt, font=\footnotesize] at ($ (P)!0.5!(E) $) {$\frac{\sqrt{3}}{3}$};

    % 底辺 BE の分割寸法 (引き出し線付き)
    \draw[thin, gray!70] ($ (B)!0.5!(Q) $) -- ($ ($ (B)!0.5!(Q) $) + (0, -0.22) $);
    \node[below=1pt, font=\footnotesize] at ($ ($ (B)!0.5!(Q) $) + (0, -0.22) $) {$\frac{2}{3}\sqrt{3}$};

    \draw[thin, gray!70] ($ (Q)!0.5!(E) $) -- ($ ($ (Q)!0.5!(E) $) + (0, -0.22) $);
    \node[below=1pt, font=\footnotesize] at ($ ($ (Q)!0.5!(E) $) + (0, -0.22) $) {$\frac{\sqrt{3}}{3}$};

  \end{tikzpicture}
  \caption{平面 $ABE$ による切断面と各垂線および寸法}
  \label{fig:cross_section_abe_sketch_reproduced}
\end{figure}

$O$ は $\triangle ACD$ の外心$P$を通り $\triangle ACD$ に垂直な直線上かつ $\triangle BCD$ の外心$Q$を通り $\triangle BCD$ に垂直な直線上にある．$\triangle ACD,BCD$は正三角形だから，$\triangle ABE$上にある．よって$O$もこの$\triangle ABE$上にある．さらに$O$から$AB$におろした垂足を$R$とする．

% TODO: triangle ACD, E, P
\begin{figure}[htb]
  \centering
  \begin{tikzpicture}[scale=2.4, >=stealth, every node/.style={font=\small}]

    % --- 1. 定数・座標パラメータ (一辺 2 の正三角形) ---
    % 底辺 CD: 長さ 2 (C(-1, 0), D(1, 0), E(0, 0))
    % 頂点 A(0, sqrt(3)) ≈ (0, 1.732)
    % 重心 P(0, sqrt(3)/3) ≈ (0, 0.577)
    \coordinate (E) at (0, 0);
    \coordinate (C) at (-1.0, 0);
    \coordinate (D) at (1.0, 0);
    \coordinate (A) at (0, {sqrt(3.0)});
    \coordinate (P) at (0, {sqrt(3.0)/3.0});

    % --- 2. 正三角形 ACD の塗りつぶしと輪郭線 ---
    \fill[blue!10, opacity=0.7] (C) -- (D) -- (A) -- cycle;
    \draw[very thick, blue!85!black] (C) -- (D) -- (A) -- cycle;

    % 中線 AE (高さ: AE ⊥ CD)
    \draw[very thick, teal!85!black] (A) -- (E);

    % 外心 P から各頂点への半径線 (破線)
    \draw[dashed, gray!60] (P) -- (C);
    \draw[dashed, gray!60] (P) -- (D);

    % --- 3. 直角記号 (E における AE ⊥ CD) ---
    \draw[thin, gray!70] (0.08, 0) -- (0.08, 0.08) -- (0, 0.08);

    % --- 4. 頂点プロットとラベル ---
    \node[circle, fill, inner sep=1.0pt] at (A) {}; \node[above=2pt] at (A) {$A$};
    \node[circle, fill, inner sep=1.0pt] at (C) {}; \node[below left=1pt] at (C) {$C$};
    \node[circle, fill, inner sep=1.0pt] at (D) {}; \node[below right=1pt] at (D) {$D$};
    \node[circle, fill=gray!80!black, inner sep=0.8pt] at (E) {}; \node[gray!80!black, below=2pt] at (E) {$E$（中点）};
    \node[circle, fill=red!85!black, inner sep=1.2pt] at (P) {}; \node[right=3pt, red!90!black, font=\footnotesize\bfseries] at (P) {$P$};

    % --- 5. 寸法ラベル ---
    % 辺の長さ 2
    \node[above left=1pt, font=\scriptsize]  at ($ (A)!0.5!(C) $) {$2$};
    \node[above right=1pt, font=\scriptsize] at ($ (A)!0.5!(D) $) {$2$};
    \node[below=2pt, font=\scriptsize]       at ($ (C)!0.5!(E) $) {$1$};
    \node[below=2pt, font=\scriptsize]       at ($ (E)!0.5!(D) $) {$1$};

    % 中線 AE の分割比 (2 : 1)
    \draw[<->, thick, purple!85!black] (-0.08, {sqrt(3.0)/3.0}) -- (-0.08, {sqrt(3.0)})
    node[midway, left=1pt, font=\scriptsize] {$|AP| = \frac{2\sqrt{3}}{3}$};
    \draw[<->, thick, purple!85!black] (-0.08, 0) -- (-0.08, {sqrt(3.0)/3.0})
    node[midway, left=1pt, font=\scriptsize] {$|PE| = \frac{\sqrt{3}}{3}$};

  \end{tikzpicture}
  \caption{正三角形 $\triangle ACD$ の中線 $AE$ と外心$P$}
  \label{fig:equilateral_triangle_acd_circumcenter}
\end{figure}

まず$P,Q$が正三角形の外心だから上図より
\begin{align}
  |AP| = |BQ| = \dfrac{1}{\cos\dfrac{\pi}{6}} = \dfrac{2}{3}\sqrt{3} \label{eq:1}
\end{align}
であり，対応して
\begin{align}
  |PE| = |QE| = |AE|-|AP| = \sqrt{3} - \dfrac{2}{3}\sqrt{3} = \dfrac{\sqrt{3}}{3} \label{eq:2}
\end{align}
である．

以下，$|ER|$をふた通りで表すことで$r$に関する関係式を導く．まず$\triangle AER$に三平方の定理を用いて
\begin{align}
  |ER| = \sqrt{|AE|^2-|AR|^2} = \sqrt{(\sqrt{3})^2-\left(\dfrac{\sqrt{3}}{2}\right)^2} = \dfrac{3}{2} \label{eq:3}
\end{align}
である．一方で$\triangle ARO, \triangle OAP, \triangle OPE$に三平方の定理を用いると
\begin{align}
  |OR|
  &= \sqrt{|OA|^2 - |AR|^2} \\
  &= \sqrt{r^2 - \left(\dfrac{3}{2}\right)^2} \\
  &= \sqrt{r^2 - \dfrac{3}{4}} \label{eq:4}
\end{align}
\begin{align}
  |OP|^2 = |OA|^2 - |AP|^2 = r^2 - \left(\dfrac{2\sqrt{3}}{3}\right)^2 = r^2 - \dfrac{4}{3} \label{eq:5}
\end{align}
\begin{align}
  |OE|
  &= \sqrt{|OP|^2+|PE|^2} \\
  & = \sqrt{\left(r^2 - \dfrac{4}{3}\right)+\left(\dfrac{\sqrt{3}}{3}\right)^2} \quad (\because\text{\cref{eq:5}})\\
  &= \sqrt{r^2 - 1} \label{eq:6}
\end{align}
である．\cref{eq:3,eq:4,eq:5}から
\begin{align}
  \dfrac{3}{2} &= \sqrt{r^2 - \dfrac{3}{4}} + \sqrt{r^2 - 1} \label{eq:7}
\end{align}
である．
$t = r^2$ として，両辺正から2乗して
\begin{align}
  \dfrac{9}{4} &= 2t - \dfrac{7}{4} + 2\sqrt{\left(t - \dfrac{3}{4}\right)(t - 1)} \\
  2 - t &= \sqrt{\left(t - \dfrac{3}{4}\right)(t - 1)}
\end{align}
$t < 2$ のもとで2乗して
\begin{align}
  t^2-4t+4 &= t^2-\dfrac{7}{4}t+\dfrac{3}{4} \\
  -4t + 4 &= -\dfrac{7}{4}t + \dfrac{3}{4} \\
  \therefore
  t &= \dfrac{13}{9}
\end{align}
である．これは$t<2$を満たす．$r > 0$ から
\begin{align}
  r &= \sqrt{t} = \dfrac{\sqrt{13}}{3}
\end{align}
である．

\bigskip

\end{document}
